\documentclass[10pt]{article}
\usepackage[margin=0.85in]{geometry}
\usepackage{booktabs}
\usepackage{hyperref}
\usepackage{xcolor}
\usepackage{float}
\usepackage{caption}
\usepackage{array}
\usepackage{amsmath}
\usepackage{tikz}
\hypersetup{colorlinks=true,linkcolor=blue!55!black,citecolor=blue!55!black,urlcolor=blue!55!black}
\definecolor{sizetwo}{HTML}{9A2E22}
\definecolor{sizetwofive}{HTML}{7A5A26}
\definecolor{sizethree}{HTML}{16408F}
\definecolor{noise}{HTML}{DDE2EC}
\captionsetup{font=small,labelfont=bf}
\input{generated/macros.tex}

\title{Open Pot: Preflop Raise Sizes in Pot-Limit Omaha Tournaments\\
When Results Are Scored by ICM}
\author{Nuttakit Kundum}
\date{October 2026}

\begin{document}
\maketitle

\begin{abstract}
A solver tree needs one answer to how big a preflop raise is, and every extra size multiplies the
tree. We ask which first-in open size, and which 3-bet size, earns the most in a six-handed
pot-limit Omaha tournament once each hand is scored in prize money by the Independent Chip Model
(ICM) rather than in chips. At the bubble of a real 374-entry, 51-paid event we solve the same
table once per size (2x, 2.5x, 3x and a pot-sized open of 3.5 big blinds), with the full game played
to the river, and measure hand by hand what each size earns the player who chooses it. With equal
30-big-blind stacks the four sizes are within run-to-run noise of each other: the largest
difference at any seat is \EqLargest\ units of prize equity (about \$\EqLargestDollars). With
uneven stacks the pot-sized open earns the most in all four tables we built, a deep or a 12-big-blind
opener, under the gun or on the button, by $0.025$ to $0.085$ units, two to ten times the gap
between independent runs. Half-pot and three-quarter-pot 3-bets are not better than a pot 3-bet on
average. The mechanism is pot-limit itself: no stack above 3.5 big blinds can open all-in, so the
open size works only through how cheaply the other stacks can call, and the pot open makes calling
most expensive. We therefore keep a single pot-sized open and pot-sized 3-bet for later work.
\end{abstract}

\section{Question}
No-limit Hold'em tournaments settled on small opening raises long ago. Pot-limit Omaha (PLO) is a
different game. The largest legal first raise is a pot raise, preflop equities run close, and with
four hole cards most of a hand's value is decided after the flop. The no-limit answer cannot simply
be copied.

The question matters for a practical reason. We plan to train a neural network on hundreds of solved
tournament spots, and each extra raise size makes every solve slower and every tree larger. Before
building that library we want the smallest set of sizes that loses nothing. That splits into two
questions: which size should the first player in use, and when a player 3-bets, should the raise be
pot-sized or smaller?

\section{Setup}
\subsection{Tournament and payouts}
The payouts are the real GGPoker ``PLO Monster Stack \$25'' structure: 374 entries, an \$8{,}602
prize pool and 51 paid places, from \$1{,}531.94 for first to \$55 for 35th to 51st. Every spot is
placed at 60 players left, nine from the money, with $60 \times 30 = 1{,}800$ big blinds in play.
Every player posts a 0.12 big-blind ante, as in hand histories from the same event. Following GG's
rule, antes are excluded from the preflop pot limit, so a pot-sized open is 3.5 big blinds and a
pot 3-bet over it reaches about 12.

\subsection{Scoring in prize money}
As in earlier work on tournament equilibria \cite{ganzfried}, payoffs are prize equity rather than
chips. Each terminal outcome of a hand is scored by Malmuth--Harville ICM \cite{harville} over the
whole field: the six players at the table, plus the 54 players at other tables, who are treated as one
crowd sharing the remaining chips at their average stack. Because the crowd's members are
interchangeable, the recursion tracks only how many of them have been placed, which keeps the
computation exact. Every outcome of every terminal node, chopped pots included, is priced once before
solving (about eight million outcomes in about 25 seconds) and looked up during training. With equal
30-big-blind stacks the bubble factor between any two players is 1.58, so calling an all-in without
pot odds needs about 61\% equity instead of 50\%. Values are reported in chips of prize equity:
prizes are scaled to the chips in play, so one unit is \$\DollarsPerUnit\ and $0.1$ units is roughly
a tenth of a big blind.

\subsection{Solver}
We use external-sampling Monte Carlo counterfactual regret minimisation \cite{lanctot} with linear
discounting over the full game: pot-sized bets and raises on every street, up to four raises preflop
and one per street afterwards, exact side pots and Omaha's two-plus-three showdown rule. Preflop hands
are grouped into 780 buckets; the flop, turn and river use 120, 120 and 10 buckets. Six-player
MCCFR carries no convergence guarantee, but it is the method behind the strongest multiplayer poker
agents \cite{pluribus}.

The engine allows one raise size per decision. Instead of offering a menu of sizes, we solve
\emph{the same table once per size}, changing only the preflop rule through a subclass of the game
state. The production engine is not modified. A 2x open builds a larger tree than a pot open
(\NodesTwoX\ million against \NodesTwo\ million nodes for equal stacks), so every run receives the
same 25 million deals rather than the same time.

\subsection{How a size is scored}
A size is judged by what it earns the player who chooses it, at the decision where it is chosen. For
the open, this is each seat's value at its first-in decision, with every earlier player folded. For
the 3-bet, it is the value of the player facing a single open with everyone between folded. We draw
6{,}000 random hands for that player and deal each one 64 times. Every other player follows the
solved average strategy, and their earlier actions weight each deal by how likely they were with the
cards they were dealt. The player's own action mix comes from the solved strategy for their bucket.

\subsection{Noise}
Every arm runs twice with different random seeds. A difference counts only when it is larger than the
gap between the two runs. The sampling error of each value is between 0.005 and 0.011 units.

\section{Equal stacks: size barely matters}
All six players hold 30 big blinds. Figure~\ref{fig:equal} and Table~\ref{tab:equal} compare each
size with a pot open at every seat.

\begin{figure}[H]
\centering
\input{generated/fig_equal.tex}
\caption{First-in value relative to a pot open, equal 30bb stacks at 60 left. Dots are the mean of two
runs and the faint lines join the two runs. Red: 2x. Brown: 2.5x. Blue: 3x. The black line is the pot
open.}
\label{fig:equal}
\end{figure}

\begin{table}[H]
\centering\small
\caption{Equal stacks: first-in value relative to a pot open, the largest two-run gap at that seat,
and the raise / limp frequency with a 2x and a pot open.}
\label{tab:equal}
\input{generated/equal_table.tex}
\end{table}

Every difference is at most \EqLargest\ units. Early seats lean slightly toward pot or 3x and the
button and small blind slightly toward 2x, but no difference clearly exceeds the largest run gap at
its seat, and none is worth more than \$\EqLargestDollars\ per decision. Averaged over the five seats, 2x, 2.5x and 3x differ
from pot by \EqAvgTwo, \EqAvgTwoFive\ and \EqAvgThree\ units. The strategy changes more than the
value: a small open turns into a limping strategy. The button limps 34\% of hands with 2x and 21\%
with pot, and the small blind limps 54\% to 69\% whatever the size.

\section{Uneven stacks: the pot open wins every table}
Real tables are rarely even, and pot-limit should matter most when stack sizes differ. We built four
bubble tables: a deep 60-big-blind stack opening into 12-big-blind stacks, and a 12-big-blind stack
opening into deep ones, each from under the gun and from the button. The 3-bet stays pot-sized.

\begin{figure}[H]
\centering
\input{generated/fig_uneven.tex}
\caption{Opener value with a 2x (red) or 2.5x (brown) open, relative to a pot open, in four uneven
tables at 60 left. Bars to the left mean the pot open earned more. Grey bands show the largest
two-run gap for that table. Stacks are listed UTG / HJ / CO / BTN / SB / BB.}
\label{fig:uneven}
\end{figure}

\begin{table}[H]
\centering\small
\caption{Uneven stacks: opener value relative to a pot open, the two-run gap, the opener's raise,
limp and fold frequency, and how often the players facing the open call.}
\label{tab:uneven}
\input{generated/uneven_table.tex}
\end{table}

The pot open earns the most in all four tables, each time beyond the run gap. The largest effect is
the deep button against short blinds, where a 2x open earns \TwoATwo\ units relative to pot, about \$0.40 per
decision, against a run gap of \GapATwo. With a pot open the short blinds call 20\% of the time; facing
2x they call 46\%. The deep stack wants to isolate a short stack or take the blinds, and only the
largest legal open achieves that.

The short opener shows the same mechanism from the other side. At 12 big blinds under the gun the
solver folds 84\% to 87\% of hands whatever the size, because it can never move all-in first. When it
does play, the pot open is best and cuts the deep stacks' calls from 25\% to 11\%.

\section{3-bet size}
On the equal-stack table, with a \ThreeOpen\ open, we compared 3-bets of half and three-quarters of a
pot raise with a pot 3-bet for all fifteen opener--defender pairs (Table~\ref{tab:three}).

\begin{table}[H]
\centering\small
\caption{3-bet size, equal stacks: value of the player facing one open relative to a pot 3-bet,
over fifteen opener--defender pairs.}
\label{tab:three}
\input{generated/three_table.tex}
\end{table}

Averaged over the pairs, both smaller sizes are within $0.005$ units of pot. One pair showed a
smaller size ahead by more than the run gap: \ThreeHalfBeyond, a half-pot 3-bet. Against every
under-the-gun open, pot was best for every defender. A pot 3-bet is sufficient.

\section{Why pot-limit changes the picture}
In a no-limit tournament a 12-big-blind stack moves all-in first, and the big stacks must fold under
ICM pressure. That fold equity is the short stack's main weapon, and small opens from deeper stacks
are cheap because the all-in reply has a fixed price. Pot-limit removes this.
\begin{itemize}
\item A first-in raise is capped at 3.5 big blinds, so only stacks of 3.5 or less can open all-in.
A pot 3-bet reaches about 12 big blinds and a pot 4-bet about 37, which is all-in at 30.
\item Stacks of 15 to 30 big blinds can never jam first in. They raise, and the next player decides
at a price the opener chose.
\item Preflop equities run close, with most matchups near 60/40, so a small open is called wide and
the hand is decided after the flop, where ICM pressure is applied.
\item The open size therefore acts through one lever: how cheaply the other stacks see a flop. The
pot open is the most expensive legal price. The solver uses it to stop limping (about 2\% limps with a
pot open against 11\% to 34\% with 2x) and to roughly halve the calls it faces.
\item With equal stacks every pair of players faces the same ICM pressure and the same stack-to-pot
ratio, and the lever is close to neutral. Uneven stacks create someone to isolate or someone to keep
out, and the larger price pays.
\end{itemize}

\section{Conclusion}
For six-handed PLO tournaments scored by ICM, a single pot-sized open and a pot-sized 3-bet are the
right preflop sizes for our solver library. Choosing pot everywhere costs at most about
\EqLargest\ units at equal stacks and is the best size in every uneven table we tested. This also
means the existing engine, which already uses pot-sized raises, needs no change.

\section{Limitations}
\begin{itemize}
\item One stage only: 60 left, just before the bubble. Early in the money and the final table are untested.
\item Each solve allowed one raise size per decision. A tree offering 2x and pot together could mix the
two and might find a small gain for some hands that this design cannot see.
\item Hands are grouped into 780 preflop buckets, so hands within a bucket share one strategy.
\item The 3-bet comparison used equal stacks only.
\item A chip-EV control ran one seed per size at equal stacks and showed no consistent pattern, so we
cannot yet separate the effect of ICM from the effect of stack depth in the uneven tables.
\item Two seeds agreeing is evidence of stability, not proof of equilibrium.
\item ICM is itself a model; field data suggests it slightly undervalues big stacks \cite{kim}.
\end{itemize}

\section*{Reproducibility}
The experiment is \texttt{research/plo\_sizing/sizing.py} in the poker-harness repository; this paper's
numbers, tables and figures are generated by \texttt{research/plo\_open\_size\_icm/build\_paper.py}
from the saved run summaries. \Solves\ solves were run on one computer, about 9 to 19 minutes
each.

\begin{thebibliography}{9}
\bibitem{pluribus} N.~Brown and T.~Sandholm. Superhuman AI for multiplayer poker. \emph{Science}, 365(6456):885--890, 2019.
\bibitem{ganzfried} S.~Ganzfried and T.~Sandholm. Computing equilibria in multiplayer stochastic games of imperfect information. In \emph{IJCAI}, 2009.
\bibitem{harville} D.~A.~Harville. Assigning probabilities to the outcomes of multi-entry competitions. \emph{Journal of the American Statistical Association}, 68(342):312--316, 1973.
\bibitem{kim} J.~Kim. Empirical validation of the Independent Chip Model. In \emph{IEEE Conference on Games}, 2025.
\bibitem{lanctot} M.~Lanctot, K.~Waugh, M.~Zinkevich and M.~Bowling. Monte Carlo sampling for regret minimization in extensive games. In \emph{NeurIPS}, 2009.
\end{thebibliography}

\end{document}
